%% =============================================================================
%% TELEQUM Portfolio — Comprehensive Technical Documentation
%% =============================================================================
%% A research-grade portfolio describing the architecture, features, and
%% capabilities of the TELEQUM v2.1 Quantum-Native 6G Digital Twin Platform.
%% =============================================================================

\documentclass[11pt,a4paper]{article}

%% ─── Packages ────────────────────────────────────────────────────
\usepackage[utf8]{inputenc}
\usepackage[T1]{fontenc}
\usepackage{lmodern}
\usepackage[margin=2.4cm]{geometry}
\usepackage{graphicx}
\usepackage{booktabs}
\usepackage{longtable}
\usepackage{tabularx}
\usepackage{amsmath, amssymb}
\usepackage{enumitem}
\usepackage{hyperref}
\usepackage{xcolor}
\usepackage{listings}
\usepackage{fancyhdr}
\usepackage{titlesec}
\usepackage{tcolorbox}
\usepackage{tikz}
\usepackage{forest}
\usepackage{multicol}

%% ─── Colour Palette ─────────────────────────────────────────────
\definecolor{telequm}{RGB}{0,120,215}
\definecolor{accent}{RGB}{255,107,0}
\definecolor{quantum}{RGB}{130,60,200}
\definecolor{codebg}{RGB}{245,245,250}
\definecolor{linkblue}{RGB}{0,90,180}

\hypersetup{colorlinks=true,linkcolor=telequm,urlcolor=linkblue,citecolor=quantum}

%% ─── Listings Style ─────────────────────────────────────────────
\lstset{
  basicstyle=\ttfamily\footnotesize,
  backgroundcolor=\color{codebg},
  frame=single,
  rulecolor=\color{telequm!30},
  breaklines=true,
  columns=fullflexible,
  keywordstyle=\color{telequm}\bfseries,
  commentstyle=\color{gray},
  stringstyle=\color{accent},
  showstringspaces=false,
}

%% ─── Section Formatting ─────────────────────────────────────────
\titleformat{\section}{\Large\bfseries\color{telequm}}{\thesection}{1em}{}
\titleformat{\subsection}{\large\bfseries\color{telequm!80!black}}{\thesubsection}{1em}{}
\titleformat{\subsubsection}{\normalsize\bfseries\color{telequm!60!black}}{\thesubsubsection}{1em}{}

%% ─── Header / Footer ────────────────────────────────────────────
\pagestyle{fancy}
\fancyhf{}
\fancyhead[L]{\small\textcolor{telequm}{\textbf{TELEQUM} v2.1 Portfolio}}
\fancyhead[R]{\small\textcolor{gray}{\today}}
\fancyfoot[C]{\small\textcolor{gray}{\thepage}}
\renewcommand{\headrulewidth}{0.4pt}

%% ─── tcolorbox styles ───────────────────────────────────────────
\tcbuselibrary{skins,breakable}
\newtcolorbox{featurebox}[1]{
  enhanced, breakable,
  colback=telequm!5, colframe=telequm!70,
  title={\textbf{#1}}, fonttitle=\bfseries\color{white},
  coltitle=white, colbacktitle=telequm,
  boxrule=0.5pt, arc=3pt,
}
\newtcolorbox{mathbox}[1]{
  enhanced, breakable,
  colback=quantum!5, colframe=quantum!50,
  title={\textbf{#1}}, fonttitle=\bfseries\color{white},
  coltitle=white, colbacktitle=quantum,
  boxrule=0.5pt, arc=3pt,
}
\newtcolorbox{warnbox}[1]{
  enhanced, breakable,
  colback=accent!5, colframe=accent!70,
  title={\textbf{#1}}, fonttitle=\bfseries\color{white},
  coltitle=white, colbacktitle=accent,
  boxrule=0.5pt, arc=3pt,
}

%% =================================================================
\begin{document}

%% ─── Title Page ──────────────────────────────────────────────────
\begin{titlepage}
\centering
\vspace*{3cm}
{\Huge\bfseries\textcolor{telequm}{TELEQUM} \par}
\vspace{0.3cm}
{\Large\textcolor{telequm!70}{Quantum-Native 6G Digital Twin Platform}\par}
\vspace{0.8cm}
{\large Version 2.1 — Technical Portfolio\par}
\vspace{2cm}

\begin{tcolorbox}[enhanced, colback=telequm!5, colframe=telequm!40,
  width=0.85\textwidth, arc=5pt, boxrule=0.6pt]
\centering
\textbf{A research-grade, scenario-driven quantum telecom simulator}\\[4pt]
bridging raw radio physics, network logic, and quantum optimization\\
with an interactive Streamlit dashboard for education and experimentation.
\end{tcolorbox}

\vspace{2cm}
{\large\textbf{Key Capabilities}\par}
\vspace{0.5cm}
\begin{minipage}{0.85\textwidth}
\begin{itemize}[leftmargin=1.5em]
  \item \textbf{7 telecom optimization problems} with QUBO formulations
  \item \textbf{5 solvers}: Greedy, Simulated Annealing, Exact, Hybrid, Quantum (QAOA/VQE)
  \item \textbf{Full quantum circuit visibility}: ASCII diagrams, optimal parameters, convergence plots, measurement distributions
  \item \textbf{Digital Twin}: time-series simulation with live KPI evolution
  \item \textbf{Education Hub}: 8 interactive subtabs from Quantum~101 to Hardware~Benchmark
  \item \textbf{3GPP-compliant} channel models (TR~38.901 UMa)
  \item \textbf{Hardware benchmark tool}: device-aware solver time estimation
\end{itemize}
\end{minipage}

\vfill
{\small\textcolor{gray}{Generated \today}\par}
\end{titlepage}

%% ─── Table of Contents ──────────────────────────────────────────
\tableofcontents
\newpage

%% =================================================================
\section{Project Overview}
%% =================================================================

TELEQUM (Telecom Quantum Utility Module) is a \textbf{quantum-native 6G digital twin platform} designed to bridge the gap between raw radio physics, network optimization logic, and quantum computing algorithms. It serves as both a \textbf{research testbed} for quantum optimization in telecommunications and an \textbf{educational platform} for understanding the intersection of quantum computing and wireless networks.

\subsection{Design Philosophy}

\begin{enumerate}[label=\textbf{R\arabic*}]
  \item \textbf{Modular Architecture} — Each component is self-contained with explicit interfaces
  \item \textbf{Stateless Optimizers} — Solvers receive snapshots, return solutions, never mutate state
  \item \textbf{Explicit QUBO Formulation} — Every problem exposes \texttt{build\_qubo()}, \texttt{decode\_solution()}, \texttt{evaluate\_cost()}
  \item \textbf{Classical Baselines First} — Every quantum experiment has a classical comparison
  \item \textbf{Scenario-Driven Methodology} — Experiments defined by reproducible YAML configs
  \item \textbf{Read-Only Dashboard} — Dashboard never modifies the simulation engine state
\end{enumerate}

\subsection{Technology Stack}

\begin{tabularx}{\textwidth}{lX}
  \toprule
  \textbf{Layer} & \textbf{Technologies} \\
  \midrule
  Core Language & Python 3.12 \\
  Quantum SDK & Qiskit 1.x, qiskit-aer, qiskit-algorithms, qiskit-optimization \\
  Simulation & NumPy, SciPy; 3GPP TR~38.901 UMa channel models \\
  Dashboard & Streamlit, Plotly \\
  External Bridges & MATLAB Engine (optional), ns-3 socket (optional) \\
  Testing & pytest, pytest-cov, GitHub Actions CI \\
  \bottomrule
\end{tabularx}


%% =================================================================
\section{Package Structure \& File Map}
%% =================================================================

The project follows a modular layout with clear separation of concerns:

\begin{lstlisting}[language={},title={\textbf{Project Root}: \texttt{notebooks/}}]
notebooks/
|-- telequm/                    # Core library
|   |-- core/                   # Quantum primitives
|   |   |-- circuits.py         # Quantum circuit construction
|   |   |-- hamiltonians.py     # Hamiltonian builders
|   |   |-- network_snapshot.py # UniversalNetworkSnapshot class
|   |   +-- visualizations.py   # Bloch sphere, circuit plots
|   |-- algorithms/             # Quantum & classical solvers
|   |   |-- qaoa.py             # NetworkQAOA class
|   |   |-- vqe.py              # ResourceVQE class
|   |   |-- qml.py              # Quantum ML utilities
|   |   +-- hybrid/             # HybridSolver (quantum-first, ensemble)
|   |       +-- hybrid_solver.py
|   |-- problems/               # QUBO problem library
|   |   |-- base_problem.py     # BaseProblem ABC
|   |   |-- prb_allocation.py   # PRBAllocationProblem
|   |   +-- telecom_problems.py # 6 more: Routing, Beam, Energy,
|   |                             Handover, BSPlacement, QuantumRouting
|   |-- simulator/              # Simulation engine
|   |   |-- engine.py           # SimulationEngine (time-step loop)
|   |   |-- network_env.py      # NetworkEnvironment
|   |   |-- optimization_bridge.py # QUBO<->solver bridge
|   |   |-- event_queue.py      # Discrete-event queue
|   |   |-- traffic_models.py   # Poisson, video, IoT traffic
|   |   +-- mobility_models.py  # Pedestrian, RWP, vehicular
|   |-- telecom/                # 3GPP utilities
|   |   |-- resource_allocation.py
|   |   |-- beamforming.py
|   |   +-- network_optimization.py
|   +-- bridges/                # External tool bridges
|       |-- matlab_bridge.py
|       +-- ns3_bridge.py
|-- dashboard/                  # Streamlit web dashboard
|   |-- app.py                  # Main entry point
|   |-- components/
|   |   |-- use_case_lab.py     # Single-shot QUBO experiments
|   |   |-- digital_twin.py    # Time-series simulation viewer
|   |   |-- education_hub.py   # 8-subtab learning platform
|   |   +-- hardware_hub.py    # Hardware management
|   +-- utils/
|       |-- resource_monitor.py # CPU/RAM tracking + HW benchmark
|       |-- snapshot_manager.py # Isolated env copies + problem factory
|       |-- scenario_loader.py  # YAML/slider config builder
|       +-- plot_helpers.py     # Plotly chart functions
|-- experiments/                # Experiment infrastructure
|-- benchmarks/                 # Standardized test topologies
+-- tests/                      # pytest test suite
\end{lstlisting}


%% =================================================================
\section{Dashboard Architecture}
%% =================================================================

The Streamlit dashboard provides four main tabs, each serving a distinct purpose:

\begin{center}
\begin{tabularx}{0.95\textwidth}{lX}
  \toprule
  \textbf{Tab} & \textbf{Purpose} \\
  \midrule
  \textbf{Use-Case Lab} & Single-shot QUBO problem solving with classical vs.\ quantum comparison \\
  \textbf{Digital Twin} & Multi-timestep network simulation with live KPI evolution \\
  \textbf{Education Hub} & 8-subtab interactive learning platform \\
  \textbf{Hardware Hub} & Hardware management and configuration \\
  \bottomrule
\end{tabularx}
\end{center}


%% =================================================================
\section{Use-Case Lab}
\label{sec:usecase}
%% =================================================================

The Use-Case Lab enables \textbf{single-shot QUBO experiments} with instant classical-vs-quantum comparison.

\subsection{Workflow}

\begin{enumerate}
  \item \textbf{Network Configuration} — Sliders for BS count (2--8), UE count (2--20), area size, Tx power, frequency, seed; or load from preset/YAML.
  \item \textbf{Problem Selection} — Choose from 7 telecom problems (see Section~\ref{sec:problems}).
  \item \textbf{Solver Selection} — Classical (Greedy, SA, Exact) or Hybrid (Quantum-First, Ensemble).
  \item \textbf{Quantum Toggle} — Enable QAOA comparison with adjustable max variables (4--100).
  \item \textbf{Experiment Mode} — Single-Shot QUBO or Full Simulation with time-series.
\end{enumerate}

\subsection{QUBO Variable Limits}

The quantum variable slider ranges from 4 to \textbf{100 variables}, with tiered resource warnings:

\begin{center}
\begin{tabular}{ccl}
  \toprule
  \textbf{Range} & \textbf{Level} & \textbf{Warning} \\
  \midrule
  4--30 & \textcolor{green!60!black}{Normal} & No warnings \\
  31--50 & \textcolor{yellow!60!black}{Info} & ``Slow on single CPU; RAM estimate shown'' \\
  51--100 & \textcolor{red!80!black}{Warning} & ``Requires multi-CPU or GPU architecture'' \\
  \bottomrule
\end{tabular}
\end{center}

\subsection{Output Panels}

After execution, the Use-Case Lab displays:

\begin{itemize}
  \item \textbf{KPI Cards}: Cells, Users, QUBO Vars, Runtime
  \item \textbf{Network State}: topology + association plot, SINR heatmap
  \item \textbf{Allocation Heatmap}: UE $\times$ Cell matrix
  \item \textbf{Classical Result}: method, cost, runtime, metrics (throughput, fairness)
  \item \textbf{Quantum Result}: cost, runtime, improvement percentage
  \item \textbf{Quantum Circuit Execution} (see Section~\ref{sec:circuit_viz})
  \item \textbf{Solution Evaluation}: problem-specific KPIs (see Section~\ref{sec:eval})
  \item \textbf{Compute Resources}: wall time, peak RAM, CPU time, complexity estimate
  \item \textbf{Export}: download results as JSON
\end{itemize}


%% =================================================================
\section{Quantum Circuit Execution Visibility}
\label{sec:circuit_viz}
%% =================================================================

When quantum solving is enabled, the dashboard displays comprehensive circuit execution details:

\subsection{Circuit Statistics}

\begin{center}
\begin{tabular}{ll}
  \toprule
  \textbf{Metric} & \textbf{Description} \\
  \midrule
  Qubits & Number of qubits = QUBO variables \\
  Depth & Circuit depth (layers of parallel gates) \\
  Gates & Total gate count (Hadamard, RZZ, Rx) \\
  Shots & Measurement shots per expectation evaluation \\
  \bottomrule
\end{tabular}
\end{center}

\subsection{Circuit Diagram}

The actual QAOA circuit is rendered as an ASCII text diagram using Qiskit's \texttt{QuantumCircuit.draw("text")}. This shows the exact gate sequence applied:

\begin{lstlisting}[title={\textbf{Example}: 4-qubit QAOA circuit (p=2)}]
     +---+                           +---+
q_0: | H |--RZZ(g0)--RZZ(g0)--Rx(b0)--| ... |--meas
     +---+                           +---+
q_1: | H |--RZZ(g0)--RZZ(g0)--Rx(b0)--| ... |--meas
     +---+                           +---+
q_2: | H |--RZZ(g0)-----------Rx(b0)--| ... |--meas
     +---+                           +---+
q_3: | H |--RZZ(g0)-----------Rx(b0)--| ... |--meas
     +---+                           +---+
\end{lstlisting}

\subsection{Optimal Parameters Table}

The optimized QAOA angles are displayed per layer:

\begin{center}
\begin{tabular}{ccc}
  \toprule
  \textbf{Layer} & $\gamma$ \textbf{(cost)} & $\beta$ \textbf{(mixer)} \\
  \midrule
  1 & 0.523691 & 1.047382 \\
  2 & 0.785028 & 0.392514 \\
  \bottomrule
\end{tabular}
\end{center}

\subsection{Convergence Plot}

An interactive Plotly chart shows the QAOA/VQE optimization convergence — expectation value versus optimizer iteration, tracking how the cost decreases during the COBYLA optimization loop.

\subsection{Measurement Distribution}

A bar chart of the top 15 measured bitstrings by count, showing the probability landscape of the quantum state. The most frequent bitstring is the optimal solution.

\subsection{Convergence Info}

\begin{itemize}
  \item \textbf{Function Evaluations}: number of times the quantum circuit was executed
  \item \textbf{Converged}: whether the scipy optimizer converged successfully
  \item \textbf{Final Cost}: final expectation value
  \item \textbf{Optimizer Message}: e.g., ``Maximum number of function evaluations exceeded''
\end{itemize}


%% =================================================================
\section{Digital Twin — Live Simulation}
%% =================================================================

The Digital Twin tab runs \textbf{multi-timestep network simulations} and visualizes the evolution of network KPIs over time.

\subsection{Configuration (Sidebar)}

\begin{itemize}
  \item \textbf{Network}: BS count, UE count, area size, seed
  \item \textbf{Problem}: 7 telecom problems (same as Use-Case Lab)
  \item \textbf{Simulation}: timesteps (10--200), traffic model (Poisson/video/IoT), mobility model (pedestrian/RWP/vehicular)
  \item \textbf{Solver}: classical method, optimization interval
  \item \textbf{Quantum}: toggle QAOA comparison (4--100 vars, same tiered warnings)
\end{itemize}

\subsection{Simulation Pipeline}

\begin{enumerate}
  \item Run full time-series simulation via \texttt{SimulationEngine}
  \item Collect per-timestep metrics (throughput, SINR, fairness, active UEs)
  \item Run QUBO problem solve on the \emph{final} network state
  \item Display all results with interactive exploration
\end{enumerate}

\subsection{Output Views}

\begin{featurebox}{Digital Twin Dashboard Panels}
\begin{itemize}
  \item \textbf{KPI Banner}: throughput, SINR, Jain fairness, active UEs, timesteps
  \item \textbf{Network State}: final topology + coverage heatmap
  \item \textbf{Performance Over Time}: throughput and fairness/SINR time-series plots
  \item \textbf{Solver Costs Over Time}: classical QUBO cost at optimization steps
  \item \textbf{QUBO Solve Results}: classical vs.\ quantum comparison on final state
  \item \textbf{Quantum Circuit Execution}: same full circuit visibility as Use-Case Lab
  \item \textbf{Compute Resources}: total time, peak RAM, CPU time
  \item \textbf{Timestep Explorer}: slider to inspect any individual timestep's KPIs
\end{itemize}
\end{featurebox}


%% =================================================================
\section{Telecom Optimization Problems}
\label{sec:problems}
%% =================================================================

TELEQUM implements \textbf{7 telecom optimization problems}, each with a complete QUBO formulation, classical solver support, quantum solver support, and problem-specific evaluation metrics.

%% ─── Problem 1 ──────────────────────────────────────────────────
\subsection{PRB Allocation (Resource Scheduling)}

\begin{mathbox}{Mathematical Formulation}
\textbf{Variables:} $x_{u,b} \in \{0,1\}$ — user $u$ served by BS $b$.

\textbf{Objective:}
$$\min_{x} \sum_{u,b} -\text{SINR}_{u,b} \cdot x_{u,b} + \lambda_1 \sum_u \left(\sum_b x_{u,b} - 1\right)^2 + \lambda_2 \sum_b \text{cap\_penalty}(b)$$

\textbf{Constraints:}
\begin{itemize}[nosep]
  \item Each user served by exactly one BS: $\sum_b x_{u,b} = 1 \quad \forall u$
  \item BS capacity: $\sum_u x_{u,b} \leq C_b \quad \forall b$
\end{itemize}

\textbf{QUBO size:} $n = U \times B$ variables.
\end{mathbox}

\begin{itemize}
  \item \textbf{Classical techniques}: Proportional Fair (PF), Round Robin (RR), Max-Rate, WFQ, Hungarian Algorithm, ILP
  \item \textbf{Bottleneck}: NP-hard for joint multi-cell scheduling; real-time 1ms TTI constraint
  \item \textbf{Evaluation}: avg throughput (Mbps), sum throughput, Jain fairness index
  \item \textbf{Files}: \texttt{problems/prb\_allocation.py}, \texttt{telecom/resource\_allocation.py}
\end{itemize}

%% ─── Problem 2 ──────────────────────────────────────────────────
\subsection{Network Routing Optimization}

\begin{mathbox}{Mathematical Formulation}
\textbf{Variables:} $x_{i,j} \in \{0,1\}$ — edge $(i,j)$ included in path.

\textbf{Objective:}
$$\min_{x} \sum_{(i,j)} -\overline{\text{SINR}}_j \cdot x_{i,j}$$

\textbf{QUBO size:} $n = N^2$ variables (where $N$ = number of cells/nodes).
\end{mathbox}

\begin{itemize}
  \item \textbf{Classical}: Dijkstra/OSPF, Bellman-Ford/BGP, Multi-Commodity Flow, Segment Routing
  \item \textbf{Bottleneck}: NP-hard for multi-commodity; SDN controller latency
  \item \textbf{Evaluation}: number of edges, total path cost
  \item \textbf{File}: \texttt{problems/telecom\_problems.py} (\texttt{RoutingOptimization})
\end{itemize}

%% ─── Problem 3 ──────────────────────────────────────────────────
\subsection{Beam Selection (Beamforming Optimization)}

\begin{mathbox}{Mathematical Formulation}
\textbf{Variables:} $x_{u,b} \in \{0,1\}$ — user $u$ uses beam $b$ from codebook.

\textbf{Objective:}
$$\min_{x} \sum_{u,b} -G_{u,b} \cdot x_{u,b} + \lambda \sum_u \left(\sum_b x_{u,b} - 1\right)^2$$

\textbf{Constraint:} One beam per user. \quad \textbf{QUBO size:} $n = U \times B_{\text{codebook}}$.
\end{mathbox}

\begin{itemize}
  \item \textbf{Classical}: Zero-Forcing (ZF), MMSE, HHL quantum linear algebra
  \item \textbf{Bottleneck}: $O(N^3)$ matrix inversion scales poorly with massive MIMO
  \item \textbf{File}: \texttt{problems/telecom\_problems.py} (\texttt{BeamSelection})
\end{itemize}

%% ─── Problem 4 ──────────────────────────────────────────────────
\subsection{Energy Efficiency (Cell On/Off)}

\begin{mathbox}{Mathematical Formulation}
\textbf{Variables:} $y_c \in \{0,1\}$ — cell $c$ is ON; $x_{u,c}$ — user $u$ assigned to cell $c$.

\textbf{Objective:}
$$\min \sum_c P_c \cdot y_c - \sum_{u,c} \text{SINR}_{u,c} \cdot x_{u,c}$$

\textbf{QUBO size:} $n = N_{\text{cell}} + U \times N_{\text{cell}}$ variables.
\end{mathbox}

\begin{itemize}
  \item \textbf{Classical}: GSMA energy frameworks, column generation, threshold-based heuristics
  \item \textbf{Goal}: Net-zero 2050; 15--30\% energy savings in dense urban deployments
  \item \textbf{Evaluation}: active cells, energy saved (\%)
  \item \textbf{File}: \texttt{problems/telecom\_problems.py} (\texttt{EnergyEfficiency})
\end{itemize}

%% ─── Problem 5 ──────────────────────────────────────────────────
\subsection{Handover Optimization}

\begin{mathbox}{Mathematical Formulation}
\textbf{Variables:} $x_{u,c}$ — user $u$ assigned to cell $c$.

\textbf{Objective:}
$$\min_{x} \sum_{u,c} -\text{SINR}_{u,c} \cdot x_{u,c} + h \sum_{u} \sum_{c \neq s_u} x_{u,c}$$
where $s_u$ is the current serving cell and $h$ is the handover penalty.

\textbf{QUBO size:} $n = U \times N_{\text{cell}}$.
\end{mathbox}

\begin{itemize}
  \item \textbf{Classical}: A3/A5 event triggers, Conditional Handover (CHO, 3GPP Rel-16)
  \item \textbf{Bottleneck}: Ping-pong handovers in dense small-cell deployments
  \item \textbf{Evaluation}: number of handovers
  \item \textbf{File}: \texttt{problems/telecom\_problems.py} (\texttt{HandoverOptimization})
\end{itemize}

%% ─── Problem 6 ──────────────────────────────────────────────────
\subsection{Base Station Placement (Facility Location)}

\begin{mathbox}{Mathematical Formulation}
\textbf{Variables:} $y_k \in \{0,1\}$ — candidate site $k$ is built.

\textbf{Objective:}
$$\min_{y} \sum_k \left(c_k - \sum_u \max(\text{SINR}_{u,k}, 0)\right) y_k + \lambda \left(\sum_k y_k - B\right)^2$$

\textbf{Budget constraint:} $\sum_k y_k \leq B$. \quad \textbf{QUBO size:} $n = K$ (candidate sites).
\end{mathbox}

\begin{itemize}
  \item \textbf{Classical}: MILP, Genetic Algorithm, Simulated Annealing, radio planning tools (Atoll, Planet)
  \item \textbf{Evaluation}: sites selected, budget utilization (\%)
  \item \textbf{File}: \texttt{problems/telecom\_problems.py} (\texttt{BSPlacementProblem})
\end{itemize}

%% ─── Problem 7 ──────────────────────────────────────────────────
\subsection{Quantum Network Routing (Future Internet)}

\begin{mathbox}{Mathematical Formulation}
\textbf{Variables:} $x_{i,j} \in \{0,1\}$ — link $(i,j)$ selected for entanglement distribution.

\textbf{Objective} (fidelity maximization via additive $-\log F$):
$$\min_{x} \sum_{(i,j)} -\log F_{i,j} \cdot x_{i,j}$$
where $F_{i,j} = F_{\text{base}} \cdot e^{-d_{i,j}/2000}$ is the link fidelity.

\textbf{QUBO size:} $n = N^2$ (all possible links).
\end{mathbox}

\begin{itemize}
  \item \textbf{Context}: Quantum Key Distribution (QKD), entanglement swapping, quantum repeater networks
  \item \textbf{Evaluation}: number of hops, end-to-end fidelity
  \item \textbf{File}: \texttt{problems/telecom\_problems.py} (\texttt{QuantumNetworkRouting})
\end{itemize}


%% =================================================================
\section{Problem-Specific Evaluation Metrics}
\label{sec:eval}
%% =================================================================

After every solve, TELEQUM displays \textbf{problem-specific KPIs} comparing classical (C) and quantum (Q) solutions side-by-side:

\begin{center}
\begin{tabularx}{\textwidth}{lX}
  \toprule
  \textbf{Problem} & \textbf{Evaluation Metrics} \\
  \midrule
  PRB Allocation & C/Q: Throughput (Mbps), Jain Fairness Index \\
  Routing & C/Q: Edges Used, QUBO Cost \\
  Beam Selection & C/Q: Users Assigned, QUBO Cost \\
  Energy Efficiency & C/Q: Active Cells, Energy Saved (\%) \\
  Handover & C/Q: Number of Handovers, QUBO Cost \\
  BS Placement & C/Q: Sites Selected, Budget Utilization (\%) \\
  Quantum Routing & C/Q: Hops, End-to-End Fidelity \\
  \bottomrule
\end{tabularx}
\end{center}


%% =================================================================
\section{Solver Architecture}
%% =================================================================

\subsection{Solver Pipeline}

The optimization pipeline follows a strict stateless pattern:

\begin{enumerate}
  \item \textbf{Snapshot}: \texttt{NetworkEnvironment} $\to$ \texttt{UniversalNetworkSnapshot}
  \item \textbf{QUBO Build}: \texttt{BaseProblem.to\_qubo()} $\to$ $(Q, \text{offset}, \text{metadata})$
  \item \textbf{Ising Convert}: \texttt{OptimizationBridge.\_qubo\_to\_ising()} via $x_i = (1-Z_i)/2$
  \item \textbf{Solve}: classical baseline \textbf{and/or} quantum algorithm
  \item \textbf{Decode}: \texttt{problem.decode\_solution(x, metadata)} $\to$ allocation dict
  \item \textbf{Evaluate}: \texttt{problem.compute\_metrics()} $\to$ KPIs
\end{enumerate}

\subsection{Classical Solvers}

\begin{center}
\begin{tabularx}{\textwidth}{lXcc}
  \toprule
  \textbf{Solver} & \textbf{Method} & \textbf{Time} & \textbf{Space} \\
  \midrule
  Greedy & Variable-by-variable QUBO minimization & $O(n^3)$ & $O(n^2)$ \\
  Simulated Annealing & Flip-based SA with exponential cooling & $O(n^2 T)$ & $O(n^2)$ \\
  Exact Brute-Force & Enumerate all $2^n$ bitstrings & $O(2^n n)$ & $O(n)$ \\
  \bottomrule
\end{tabularx}
\end{center}

\subsection{Quantum Solvers}

\begin{center}
\begin{tabularx}{\textwidth}{lXc}
  \toprule
  \textbf{Algorithm} & \textbf{Description} & \textbf{File} \\
  \midrule
  \textbf{QAOA} & Quantum Approximate Optimization Algorithm; $p$-layer parameterized circuit with cost ($e^{-i\gamma H_C}$) and mixer ($e^{-i\beta H_M}$) unitaries. Optimized via COBYLA. & \texttt{qaoa.py} \\
  \textbf{VQE} & Variational Quantum Eigensolver; RealAmplitudes ansatz with Ry rotations and CNOT entangling layers. & \texttt{vqe.py} \\
  \bottomrule
\end{tabularx}
\end{center}

\subsection{Hybrid Solvers}

The \texttt{HybridSolver} class (\texttt{algorithms/hybrid/hybrid\_solver.py}) implements two strategies:

\begin{itemize}
  \item \textbf{Quantum-First}: attempt quantum solve; fall back to classical if the problem exceeds qubit limits
  \item \textbf{Ensemble}: run all available solvers (greedy, SA, quantum) and select the best solution by QUBO cost
\end{itemize}

\subsection{Solver Output (with Circuit Info)}

The \texttt{OptimizationBridge.solve\_quantum()} method returns a rich result dictionary:

\begin{center}
\begin{tabularx}{\textwidth}{lX}
  \toprule
  \textbf{Field} & \textbf{Description} \\
  \midrule
  \texttt{solution} & Binary vector $x \in \{0,1\}^n$ \\
  \texttt{decoded} & Human-readable allocation (problem-specific) \\
  \texttt{cost} & QUBO cost $x^T Q x + \text{offset}$ \\
  \texttt{runtime\_s} & Wall-clock time (seconds) \\
  \texttt{circuit\_text} & ASCII circuit diagram \\
  \texttt{optimal\_params} & Final $[\gamma_0..\gamma_p, \beta_0..\beta_p]$ \\
  \texttt{optimization\_history} & Cost per optimizer iteration \\
  \texttt{measurement\_counts} & Bitstring $\to$ count dictionary \\
  \texttt{circuit\_depth} & Circuit depth \\
  \texttt{gate\_count} & Total gate count \\
  \texttt{num\_qubits} & Qubits used \\
  \texttt{convergence\_info} & nfev, success, message, final\_cost \\
  \bottomrule
\end{tabularx}
\end{center}


%% =================================================================
\section{QAOA Circuit Construction}
%% =================================================================

The QAOA circuit for a general QUBO with Ising Hamiltonian $H_C$ is constructed as:

\subsection{QUBO to Ising Conversion}

Using the substitution $x_i = \frac{1 - Z_i}{2}$, the QUBO objective $\min x^T Q x$ becomes an Ising Hamiltonian:

$$H_C = \sum_{i<j} \frac{Q_{ij}}{4} Z_i Z_j - \sum_i \left(\frac{Q_{ii}}{2} + \sum_{j>i} \frac{Q_{ij}}{4}\right) Z_i + \text{const}$$

\subsection{Circuit Layers}

For each of $p$ QAOA layers:

\begin{enumerate}
  \item \textbf{Initial State}: Apply Hadamard to all qubits: $|+\rangle^{\otimes n}$
  \item \textbf{Cost Layer} ($e^{-i\gamma_k H_C}$): For each non-zero $ZZ$ term: CNOT $\to$ Rz($\gamma$) $\to$ CNOT
  \item \textbf{Mixer Layer} ($e^{-i\beta_k H_M}$): Apply Rx($2\beta_k$) to each qubit
  \item \textbf{Measurement}: measure all qubits in computational basis
\end{enumerate}

\subsection{Gate Count Analysis}

\begin{center}
\begin{tabular}{lc}
  \toprule
  \textbf{Component} & \textbf{Gates per Layer} \\
  \midrule
  Cost layer (per edge) & 2 CNOT + 1 Rz = 3 gates \\
  Mixer layer (per qubit) & 1 Rx \\
  \midrule
  \textbf{Total per layer} & $3|E| + n$ \\
  \textbf{Total circuit} & $p(3|E| + n) + n$ (initial H gates) \\
  \bottomrule
\end{tabular}
\end{center}


%% =================================================================
\section{Education Hub — 8 Interactive Subtabs}
%% =================================================================

The Education Hub provides a comprehensive learning platform with 8 subtabs:

\subsection{Subtab 1: Quantum Computing 101}

Covers fundamental quantum concepts: qubits, superposition, entanglement, Bloch sphere visualization, quantum gates (X, Y, Z, H, CNOT, Rz), and the quantum advantage thesis.

\subsection{Subtab 2: Memory \& Complexity Calculator}

Interactive tools for:
\begin{itemize}
  \item \textbf{Statevector RAM estimation}: $2^n \times 16$ bytes (complex128)
  \item \textbf{Big-O complexity analysis}: visual scaling comparison across algorithms
  \item \textbf{Feasibility assessment}: laptop ($<16$\,GB) vs.\ server ($<256$\,GB) thresholds
\end{itemize}

\subsection{Subtab 3: Algorithm Deep Dives}

Detailed explanations of each solver, organized as:
\begin{center}
\begin{tabularx}{\textwidth}{lX}
  \toprule
  \textbf{Algorithm} & \textbf{Content} \\
  \midrule
  QAOA & Math formulation $\to$ QUBO conversion $\to$ Ising map $\to$ circuit $\to$ measurement $\to$ code mapping \\
  VQE & Ansatz design $\to$ parameter optimization $\to$ convergence analysis \\
  Greedy & Step-by-step variable selection strategy $\to$ complexity proof \\
  SA & Temperature schedule $\to$ acceptance probability $\to$ convergence \\
  Exact & Enumeration strategy $\to$ infeasibility threshold \\
  \bottomrule
\end{tabularx}
\end{center}

\subsection{Subtab 4: Quantum Circuits Explained}

Gate-by-gate breakdown of QAOA and VQE circuit construction with ASCII diagrams, covering:
\begin{itemize}
  \item Cost unitary implementation (CNOT-Rz-CNOT)
  \item Mixer unitary implementation (Rx)
  \item Parameter binding and optimization loop
  \item Measurement and bitstring decoding
\end{itemize}

\subsection{Subtab 5: Problem Statements \& Modeling}

Detailed descriptions of all 7 telecom problems, each including:
\begin{itemize}
  \item Real-world context and industry relevance
  \item Mathematical formulation (variables, objectives, constraints)
  \item Classical vendor techniques and their limitations
  \item Scalability bottlenecks (NP-hardness, real-time constraints)
  \item QUBO formulation and quantum approach
  \item Key references (3GPP standards, academic papers)
\end{itemize}

\subsection{Subtab 6: Solver Architecture \& File Map}

Visual pipeline diagram showing how data flows from \texttt{NetworkEnvironment} through the QUBO construction, solver dispatch, and result decoding. Includes a complete file table mapping each module to its purpose.

\subsection{Subtab 7: Research References}

Curated database of 9 key research papers organized by domain:

\begin{enumerate}
  \item Choi et al.\ — Quantum Approximation for Wireless Scheduling (IEEE, 2019)
  \item Kim et al.\ — Quantum Annealing for Resource Allocation (IEEE Wireless, 2021)
  \item Farhi et al.\ — QAOA (arXiv:1411.4028, 2014)
  \item Peruzzo et al.\ — VQE (Nature Comm., 2014)
  \item Harrigan et al.\ — QAOA on sycamore (Nature Physics, 2021)
  \item 3GPP TR~38.901 — Channel Models (v16.1.0)
  \item Moll et al.\ — Quantum Optimization for Telco (IEEE Network, 2023)
  \item Caleffi et al.\ — Quantum Internet Networking (IEEE Network, 2020)
  \item Kozlowski \& Wehner — Quantum Network Protocol Design (arXiv, 2019)
\end{enumerate}

\subsection{Subtab 8: Hardware Benchmark}

Interactive hardware profiling tool that:

\begin{featurebox}{Hardware Benchmark Features}
\begin{enumerate}
  \item \textbf{Detects Device Specs}: CPU name, cores, frequency; RAM total/available; GPU (Apple Metal or NVIDIA); OS; Python version; architecture
  \item \textbf{Runs Performance Benchmark}: 500$\times$500 matrix multiply to compute GFLOPS score
  \item \textbf{Estimates Solver Times}: for any QUBO size $n$, shows:
    \begin{itemize}
      \item Per-solver: Big-O notation, operation count, estimated wall-clock time, RAM, feasibility
      \item Comparison table at multiple scales ($n = 5, 10, 15, 20, 25, 30, 40, 50, 60$)
    \end{itemize}
  \item \textbf{Performance Tier Assessment}: High ($>$20 GFLOPS), Standard (5--20), Low ($<$5)
\end{enumerate}
\end{featurebox}


%% =================================================================
\section{Simulation Engine}
%% =================================================================

The \texttt{SimulationEngine} (\texttt{simulator/engine.py}) drives the Digital Twin's time-series simulation:

\subsection{Engine Loop}

\begin{enumerate}
  \item Initialize \texttt{NetworkEnvironment} from config
  \item For each timestep $t = 0, 1, \ldots, T$:
    \begin{enumerate}
      \item Update user positions via mobility model
      \item Generate new traffic arrivals
      \item Recalculate channel (path loss, SINR) using 3GPP TR~38.901
      \item At optimization intervals: take snapshot $\to$ build QUBO $\to$ solve $\to$ apply allocation
      \item Record metrics (throughput, SINR, fairness, active UEs)
    \end{enumerate}
  \item Return full metrics history + final environment state
\end{enumerate}

\subsection{Channel Model}

SINR is computed using the 3GPP TR~38.901 Urban Macro (UMa) path loss model:

$$\text{PL}[\text{dB}] = 28.0 + 22 \log_{10}(d) + 20 \log_{10}(f_c)$$

where $d$ is the 3D distance in meters and $f_c$ is the carrier frequency in GHz.

\subsection{Traffic Models}

\begin{center}
\begin{tabular}{ll}
  \toprule
  \textbf{Model} & \textbf{Characteristics} \\
  \midrule
  Poisson & Constant arrival rate $\lambda$, exponential inter-arrival times \\
  Video & Bursty traffic with periodic large packets \\
  IoT & Low-rate, sporadic transmissions \\
  \bottomrule
\end{tabular}
\end{center}

\subsection{Mobility Models}

\begin{center}
\begin{tabular}{ll}
  \toprule
  \textbf{Model} & \textbf{Characteristics} \\
  \midrule
  Pedestrian & Low speed ($\sim$1.2 m/s), random direction changes \\
  Random Waypoint & Pick random destination, move in straight line, pause \\
  Vehicular & High speed ($\sim$30 m/s), road-constrained movement \\
  \bottomrule
\end{tabular}
\end{center}


%% =================================================================
\section{Resource Monitoring}
%% =================================================================

The \texttt{resource\_monitor.py} module provides real-time resource tracking:

\subsection{Track Resources Context Manager}

\begin{lstlisting}[language=Python]
with track_resources() as report:
    result = solve_problem(...)
print(report.wall_time_s, report.peak_ram_mb, report.cpu_time_s)
\end{lstlisting}

Uses \texttt{tracemalloc} for peak RAM and \texttt{time.process\_time()} for CPU time.

\subsection{Estimation Functions}

\begin{center}
\begin{tabularx}{\textwidth}{lX}
  \toprule
  \textbf{Function} & \textbf{Purpose} \\
  \midrule
  \texttt{estimate\_statevector\_ram(n)} & RAM for $2^n$ complex128 entries \\
  \texttt{estimate\_qaoa\_resources(n, p, |E|)} & Gates, depth, parameters, RAM \\
  \texttt{estimate\_vqe\_resources(n, layers)} & Ansatz gates/depth/parameters \\
  \texttt{estimate\_classical\_resources(n, method)} & Big-O time/space per method \\
  \texttt{benchmark\_device()} & CPU/RAM/GPU specs + GFLOPS score \\
  \texttt{estimate\_solver\_times(n, hw)} & Per-solver wall-clock time estimate \\
  \bottomrule
\end{tabularx}
\end{center}


%% =================================================================
\section{External Bridges}
%% =================================================================

TELEQUM supports optional integration with external tools:

\subsection{MATLAB Bridge}

\texttt{bridges/matlab\_bridge.py} connects to the MATLAB Engine for leveraging MATLAB's Communications Toolbox, specifically for:
\begin{itemize}
  \item Advanced channel modeling beyond 3GPP baseline
  \item Proprietary beamforming algorithms
  \item Result validation against MATLAB reference implementations
\end{itemize}

\subsection{ns-3 Bridge}

\texttt{bridges/ns3\_bridge.py} connects to ns-3 via socket for:
\begin{itemize}
  \item Packet-level network simulation
  \item MAC/PHY layer validation
  \item End-to-end latency and jitter analysis
\end{itemize}


%% =================================================================
\section{Testing \& Benchmarking}
%% =================================================================

\subsection{Test Suite}

The \texttt{tests/} directory contains pytest-based tests:
\begin{itemize}
  \item \texttt{test\_simulator.py}: engine initialization, snapshot creation, QUBO construction, solver execution
\end{itemize}

Run with: \texttt{pytest tests/ -v --cov=telequm}

\subsection{Benchmarks}

The \texttt{benchmarks/} directory provides standardized topologies:
\begin{itemize}
  \item Small (3 BS, 6 UE), Medium (5 BS, 15 UE), Large (8 BS, 30 UE)
  \item \texttt{run\_benchmarks.py}: automated benchmark runner with CSV output
\end{itemize}


%% =================================================================
\section{Complete File Reference}
%% =================================================================

\begin{longtable}{lp{8cm}}
\toprule
\textbf{File} & \textbf{Purpose} \\
\midrule
\endfirsthead
\toprule
\textbf{File} & \textbf{Purpose} \\
\midrule
\endhead

\multicolumn{2}{l}{\textbf{telequm/core/}} \\
\texttt{circuits.py} & Quantum circuit construction helpers \\
\texttt{hamiltonians.py} & Hamiltonian builder (Ising, Heisenberg) \\
\texttt{network\_snapshot.py} & \texttt{UniversalNetworkSnapshot}, \texttt{CellInfo}, \texttt{UserInfo} \\
\texttt{visualizations.py} & Bloch sphere plots, circuit rendering \\
\midrule

\multicolumn{2}{l}{\textbf{telequm/algorithms/}} \\
\texttt{qaoa.py} & \texttt{NetworkQAOA}: QAOA with COBYLA optimizer \\
\texttt{vqe.py} & \texttt{ResourceVQE}: VQE with RealAmplitudes ansatz \\
\texttt{qml.py} & Quantum machine learning utilities \\
\texttt{hybrid/hybrid\_solver.py} & \texttt{HybridSolver}: quantum-first \& ensemble strategies \\
\midrule

\multicolumn{2}{l}{\textbf{telequm/problems/}} \\
\texttt{base\_problem.py} & \texttt{BaseProblem} ABC with \texttt{to\_qubo()}, \texttt{solve\_quantum()}, etc. \\
\texttt{prb\_allocation.py} & \texttt{PRBAllocationProblem} \\
\texttt{telecom\_problems.py} & 6 problems: Routing, Beam, Energy, Handover, BSPlacement, QuantumRouting \\
\midrule

\multicolumn{2}{l}{\textbf{telequm/simulator/}} \\
\texttt{engine.py} & \texttt{SimulationEngine}: time-step loop \\
\texttt{network\_env.py} & \texttt{NetworkEnvironment}: BS/UE state, SINR calculation \\
\texttt{optimization\_bridge.py} & \texttt{OptimizationBridge}: QUBO$\leftrightarrow$solver interface (with circuit info) \\
\texttt{event\_queue.py} & Priority-based discrete-event queue \\
\texttt{traffic\_models.py} & Poisson, video, IoT traffic generators \\
\texttt{mobility\_models.py} & Pedestrian, RWP, vehicular mobility \\
\midrule

\multicolumn{2}{l}{\textbf{telequm/telecom/}} \\
\texttt{resource\_allocation.py} & 3GPP-compliant resource allocation \\
\texttt{beamforming.py} & Beamforming vector computation \\
\texttt{network\_optimization.py} & Network-level optimization utilities \\
\midrule

\multicolumn{2}{l}{\textbf{telequm/bridges/}} \\
\texttt{matlab\_bridge.py} & MATLAB Engine integration \\
\texttt{ns3\_bridge.py} & ns-3 socket integration \\
\midrule

\multicolumn{2}{l}{\textbf{dashboard/}} \\
\texttt{app.py} & Streamlit main entry point \\
\texttt{components/use\_case\_lab.py} & Single-shot QUBO experiment tab \\
\texttt{components/digital\_twin.py} & Time-series simulation tab \\
\texttt{components/education\_hub.py} & 8-subtab learning platform \\
\texttt{components/hardware\_hub.py} & Hardware management tab \\
\texttt{utils/resource\_monitor.py} & Resource tracking + HW benchmark \\
\texttt{utils/snapshot\_manager.py} & Problem factory + isolated env copies \\
\texttt{utils/scenario\_loader.py} & YAML/slider config builder \\
\texttt{utils/plot\_helpers.py} & Plotly chart functions \\
\midrule

\multicolumn{2}{l}{\textbf{Other}} \\
\texttt{experiments/metrics.py} & Experiment metrics collection \\
\texttt{experiments/visualize.py} & Experiment visualization \\
\texttt{benchmarks/topologies.py} & Standardized test topologies \\
\texttt{run\_experiment.py} & CLI experiment runner \\
\texttt{run\_benchmarks.py} & CLI benchmark runner \\
\texttt{tests/test\_simulator.py} & pytest test suite \\
\bottomrule
\end{longtable}


%% =================================================================
\section{Running the Platform}
%% =================================================================

\subsection{Installation}

\begin{lstlisting}[language=bash]
cd notebooks/
python -m venv .venv
source .venv/bin/activate
pip install -r requirements.txt
\end{lstlisting}

\subsection{Launch Dashboard}

\begin{lstlisting}[language=bash]
cd notebooks/
streamlit run dashboard/app.py
\end{lstlisting}

\subsection{Run Experiments (CLI)}

\begin{lstlisting}[language=bash]
python run_experiment.py --config experiments/configs/default.yaml
python run_benchmarks.py --output results/
\end{lstlisting}

\subsection{Run Tests}

\begin{lstlisting}[language=bash]
pytest tests/ -v --cov=telequm
\end{lstlisting}


\end{document}
